\chapter{Lineaire vergelijkingen en functies}
\label{chap:functies}

In dit hoofdstuk gaan we leren over lineaire vergelijkingen en lineaire 
functies. Dat zijn niet zomaar losse onderwerpen: ze vormen samen de 
eerste echte stap van het rekenen met getallen naar het werken met 
\emph{relaties tussen grootheden}.

Een lineaire functie beschrijft een verband waarin \'e\'en grootheid 
\emph{evenredig} toeneemt met een andere. Denk aan de kosten van een 
taxirit: 
\[
    y = 4.50 + 2.20x
\]
met $x$ het aantal gereden kilometers en \EUR{4.50} het startbedrag. Of neem de temperatuuromrekening van Fahrenheit naar Celsius:
\[
    F = \frac{9}{5}C + 32
\]
of de afstand die een auto aflegt bij een constante snelheid: $s = v \cdot t$. Al deze verbanden hebben dezelfde vorm:
\[
    y = a\cdot x + b
\]

Hierin is $a$ de \emph{richtingsco\"effici\"ent} van de lijn (hoe sterk $y$ verandert als $x$ met \'e\'en toeneemt) en 
$b$ het \emph{begingetal} (de waarde van $y$ als $x = 0$). In een 
grafiek is dit een \emph{rechte lijn} (vandaar ook de naam 
\emph{lineair}, van het Latijnse \textit{linea}, ``lijn'').

Een \emph{lineaire vergelijking} ontstaat zodra je zo'n verband 
gelijk stelt aan een concrete waarde. Bijvoorbeeld: ``na hoeveel 
kilometers kost de taxirit $25$ euro?'' leidt tot de vergelijking 
$4{,}50 + 2{,}20x = 25$. Het oplossen van zulke vergelijkingen (met 
de technieken die je in de vorige hoofdstukken hebt opgedaan, zoals 
rekenen met breuken en haakjes wegwerken) staat dan ook centraal in 
dit hoofdstuk.

% In dit hoofdstuk gaan we leren over lineaire vergelijkingen en lineaire functies.

\section{Lineaire vergelijking}
Een lineaire vergelijking (of eerstegraadsvergelijking) is een vergelijking van de vorm:

\regel{}{a\cdot x+b=0}

In deze vergelijking is $x$ de onbekende, terwijl $a$ en $b$ constanten zijn. Omdat de onbekende $x$ de exponent 1 heeft, noemen we dit een lineaire vergelijking. Een lineaire vergelijking kunnen we oplossen door de termen met $x$ naar het linkerlid en de overige termen naar het rechterlid te brengen.

Voorbeelden:
\begin{enumerate}
    \itemsep1em
    \item Los op: $2x+8=0$ \\
            Oplossing: 	
            \begin{align*}
                2x+8    &=0 \\
                2x      &=-8 \\
                x       &=-4
            \end{align*} \\
	\item Los op: $x-5=4x+7$
	\item Los op: $3(x+2)+2=2(1-x)-x$ \\
        Oplossing: 	
        \begin{align*}
        3(x+2)+2  &= 2(1-x)-x \\
        3x+6+2    &= 2-2x-x  \\
        3x+8      &=2-3x \\
        3x+3x     &=2-8  \\
        6x        &=-6 \\
        x         &=-1
        \end{align*} \\
	\item Los op: $3(x-2)+2x=-4(3-x)+3x$
\end{enumerate}

\section{Lineaire functie}
Algemeen: een functie is een voorschrift om volgens een bepaalde regel uit een gegeven getal een ander getal, de functiewaarde, te maken. In dit hoofdstuk beperken we ons tot de lineaire functie.

Een lineaire functie (of $1^e$ graadsfunctie) is een functie van de vorm: 	
\regel{}{f(x)=a\cdot x+b}
In plaats van $f(x)=a\cdot x+b$ gebruiken we ook de notatie:
\regel{}{y=a\cdot x+b}

In deze functie is $x$ de (onafhankelijke) variabele, terwijl $a$ en $b$ constanten zijn. Omdat de variabele $x$ de exponent 1 heeft, noemen we dit een lineaire functie.

Een voorbeeld van een lineaire functie is: 
\[
    f(x)=2x-3
\]

De grafiek van een lineaire functie is een rechte lijn.
In $y=a\cdot x+b$ wordt $a$ de richtingsco\"effici\"ent van de rechte lijn genoemd, afgekort als rc en ``rico'' genoemd. De richtingsco\"effici\"ent $a$ geeft de helling van de lijn weer.\\

Bij een positieve $a$ hoort een stijgende grafiek, bij een negatieve $a$ een dalende grafiek.

Voorbeelden:
\begin{enumerate}
    \itemsep1em
    \item Teken de grafiek van de functie $y=2x+4$ \\
        Oplossing: Het snijpunt met de X-as vinden we uit: 
        \begin{align*}
            y=0 \Rightarrow 0&=2x+4 \\
            -2x&=4 \\
              x&=-2
        \end{align*}
        Dus: $S_x (-2,0)$ \\
        
        Opmerking: Door de lineaire functie $y=0$ te stellen, hebben we een lineaire vergelijking in $x$ gekregen! \\
        Het snijpunt met de Y-as vinden we uit: $x=0 \Rightarrow y=4$ \\
        
        Dus: $S_Y (0,4)$ \\
        \begin{figure}[H]
            \centering
            \begin{tikzpicture}[xscale=2.0, yscale=1.0]
            % Grid
            % \draw[gray!30, very thin] (-2.5,-1.5) grid (1.5,4.5);
            
            % Axes
            \draw[->,thick] (-2.5,0) -- (1.7,0) node[below] {$x$};
            \draw[->,thick] (0,-1.5) -- (0,4.7) node[left]  {$y$};
            
            % Tick marks x-axis
            \foreach \x in {-2,-1,1}
                \draw (\x,0.05) -- (\x,-0.05) node[below] {\small $\x$};
            \node[below left] at (0,0) {\small $0$};
            
            % Tick marks y-axis
            \foreach \y in {2,4}
                \draw (0.05,\y) -- (-0.05,\y) node[left] {\small $\y$};
            
            % Line: passes through (-2,0) and (0,4), slope = 2
            \draw[blue, thick] (-2.5,-1) -- (1,6);
            \end{tikzpicture}
            \caption{Grafiek van $y=2x+4$}
            \label{fig:placeholder}
        \end{figure}
    \item Teken de grafiek van de functie $y=-3x+6$
    \item Bereken de richtingsco\"effici\"ent van de lijn $x+2y=6$ en het snijpunt van deze lijn met de Y-as.
    Oplossing:
        Eerst gaan we $x+2y=6$ herschrijven in de vorm $y=ax+b$:
        \begin{align*}
            x+2y&=6 \\
            2y&=-x+6 \\
            y&=-0.5x+3
        \end{align*}
        Dus: $rc=a=-0.5$ \\
        Het snijpunt met de Y-as is $S_Y (0,3)$
    \item Bereken de richtingsco\"effici\"ent van de lijn $2x-7y=-2$ en het snijpunt van deze lijn met de Y-as.

	\item Bepaal de vergelijking van de rechte lijn $l$, die gaat door $A(1,1)$ en $B(3,5)$.
    Oplossing:
        De richtingsco\"effici\"ent $a$ kunnen we bepalen uit: 
        \[a=\frac{d_y}{d_x} =\frac{y_B-y_A}{x_B-x_A}=\frac{5-1}{3-1}=\frac{4}{2}=2\]
        De lijn $l$ heeft dus als vergelijking: $y=2x+b$ \\
        De onbekende $b$ vinden door de co\"effici\"enten van \'of $A$ \'of $B$ in te vullen. Nemen we bijvoorbeeld $A(1,1)$ dan krijgen we: 
        \begin{align*}
            1&=2\cdot 1+b \\
            b&=-1
        \end{align*}
        Dus de vergelijking van $l$ wordt: $y=2x-1$

    \item Bepaal de vergelijking van de rechte lijn $l$, die gaat door $A(4,-3)$ en $B(-3,4)$.

	\item Bepaal de vergelijking van de rechte lijn $l$, die gaat door $P(3,-1)$ en $Q(-3,-1)$.
    Oplossing:
        De richtingsco\"effici\"ent a kunnen we bepalen uit: 
        \[a=\frac{d_y}{d_x} =\frac{y_Q-y_P}{x_Q-x_P}=\frac{-1-1}{-3-3}=\frac{0}{-6}=0\]
        De lijn $l$ heeft dus als vergelijking: $y=0x+b$ \\
        $P$ invullen geeft: 
        \begin{align*}
            -1&=0\cdot 3+b \\
            b&=-1
        \end{align*}
        Dus de vergelijking van $l$ wordt: $y=-1$
    \item Bepaal de vergelijking van de rechte lijn $l$, die gaat door $P(3,-1)$ en $Q(5,3)$.
\end{enumerate}

\section{Evenwijdig en loodrecht}
Er zijn situaties waarbij voor twee rechte lijnen een bijzondere eigenschap geldt.
Twee rechte lijnen, $l_1$ en $l_2$, kunnen evenwijdig aan elkaar zijn. Notatie: $l_1 || l_2$.
Ook kunnen de lijnen $l_1$ en $l_2$ elkaar loodrecht (of: haaks) snijden. Notatie: $l_1 \perp l_2$.

Als gegeven is:	
\begin{align*}    
    l_1:y&=a_1\cdot x+b_1 \\
    l_2:y&=a_2\cdot x+b_2
\end{align*}

Dan gelden de volgende eigenschappen:
\regel{Eigenschap 1}{\text{als } l_1\ ||\ l_2 \text{, dan: }a_1=a_2}
\regel{Eigenschap 2}{\text{als } l_1\perp l_2 \text{, dan: } a_1\cdot a_2=-1}

Voorbeelden:
\begin{enumerate}
    \itemsep1em
    \item Gegeven de lijn $l_1:4x-2y=9$ en het punt $A(5,3)$. \\
          Bepaal de vergelijking van de lijn $l_2$ door $A$ en evenwijdig aan $l_1$.\\ 

          Oplossing: De vergelijking van $l_1$ is te schrijven als $y=2x-4.5$.\\
          De richtingsco\"effici\"ent van $l_1$ is dus $2$ en dit is gelijk aan de richtingsco\"effici\"ent van $l_2$. \\
          De vergelijking van $l_2$ wordt dan $y=2x+b$. \\
          Hieruit kunnen we $b$ oplossen door punt $A$ in te vullen: 
          \begin{align*}
            3&=2\cdot 5+b \\
            b&=-7 
          \end{align*}
          en dus wordt de vergelijking van: $l_2:y=2x-7$

    \item Gegeven de lijn $l_1:x+3y=5$ en het punt $A(-2,-4)$.
          Bepaal de vergelijking van de lijn $l_2$ door $A$ en loodrecht op $l_1$

    \item Gegeven zijn de lijnen $l_1:y=2x-6$ en $l_2:y=-3x+4$. \\
          Bepaal de co\"ordinaten van het snijpunt van $l_1$ en $l_2$. \\
          Oplossing: Om het snijpunt te bepalen kunnen we de vergelijkingen aan elkaar gelijkstellen, dan krijgen we: 
          \begin{align*}
            2x-6&=-3x+4 \\
            2x+3x&=4+6 \\
            5x&=10 \\
            x&=2.
          \end{align*}
            Nu kunnen we de y-co\"ordinaat vinden door de gevonden $x$ in een van de twee vergelijkingen in te vullen: 
            \begin{align*}
                y&=2\cdot 2-6 \\
                 &=4-6 \\
                 &=-2
            \end{align*}
            We hebben nu het snijpunt gevonden: $(2,-2)$

    \item Gegeven zijn de lijnen$ l_1:y=3x+7$ en $l_2:y=x-3$. \\
          Bepaal de co\"ordinaten van het snijpunt van $l_1$ en $l_2$.
\end{enumerate}

\newpage
\section{Huiswerkopdrachten}
Hieronder vind je opgaven over het oplossen van lineaire vergelijkingen en het werken met rechte lijnen in het vlak. Gebruik de rekenregels en formules uit dit hoofdstuk. Let bij elke stap op de juiste toepassing van de regels en noteer je tussenstappen duidelijk.

\subsection*{Opgave 1: Los de volgende vergelijkingen op in $\mathbb{R}$.}
($\mathbb{R}$ is de verzameling van alle re\"ele getallen)

\renewcommand{\arraystretch}{1.3}
\begin{tabular}{m{14em}@{\hspace{4em}}l}
\text{a.}\hspace{0.5em} $3 - 2x = -6$        & \text{e.}\hspace{0.5em} $0{,}5p + 1{,}5 = 1 + 2{,}5p$ \\
\text{b.}\hspace{0.5em} $x - 12 = 3 - 4x$   & \text{f.}\hspace{0.5em} $5(-2x + 3) + (2x - 5) = 4(x - 4)$ \\
\text{c.}\hspace{0.5em} $5 - (1 - x) = 3(x - 2)$ & \text{g.}\hspace{0.5em} $2(a - 3) + a = -3(4 - a) + 4a$ \\
\text{d.}\hspace{0.5em} $2 - (x + 4) = -2(x + 1) - 3$ & \text{h.}\hspace{0.5em} $2(a - 3) + a = 3(4 - a) + 6a$ \\
\end{tabular}

\subsection*{Opgave 2: Bepaal van de volgende lijnen in het vlak de richtingsco\"effici\"ent, het snijpunt met de $X$-as en het snijpunt met de $Y$-as.}
\renewcommand{\arraystretch}{1.3}
\begin{tabular}{m{14em}@{\hspace{4em}}l}
\text{a.}\hspace{0.5em} $3x - 4y = 4$   & \text{e.}\hspace{0.5em} $5y = 7$ \\
\text{b.}\hspace{0.5em} $2x = y + 7$    & \text{f.}\hspace{0.5em} $x = 3y - 2$ \\
\text{c.}\hspace{0.5em} $-4x + 2y = 3$  & \text{g.}\hspace{0.5em} $x = 2y$ \\
\text{d.}\hspace{0.5em} $-x - 5y = 1$   & \text{h.}\hspace{0.5em} $-5x + 2y = -3$ \\
\end{tabular}

\subsection*{Opgave 3: Bepaal de vergelijking van de rechte lijn door het gegeven punt met de gegeven richtingsco\"effici\"ent.}
\begin{enumerate}
\item[a.] $P(-1,-1)$ en $rc = -2$ \\
\item[b.] $P(4,0)$ en $rc = -4$ \\
\item[c.] $P(-1,-3)$ en $rc = 0$ \\
\item[d.] $P(2,-3)$ en $rc = 4$ \\
\end{enumerate}

\subsection*{Opgave 4: Loodrecht of evenwijdig.}
\begin{enumerate}
\item[a.] De grafiek van de functie $y = a \cdot x + 6$ gaat door het punt $P$ met co\"ordinaten $(-1,-2)$. Bepaal $a$. \\

\item[b.] De lijn $y = ax - 4$ is evenwijdig aan de rechte door de punten $A(2,-3)$ en $B(-3,2)$. Bepaal de vergelijkingen van beide lijnen. \\

\item[c.] Gegeven is de lijn $l_1$ met de vergelijking $2x - 5y = 12$. Bepaal de vergelijking van de rechte lijn $l_2$, die gaat door het punt $P(2,4)$ en die evenwijdig is aan $l_1$. \\

\item[d.] Gegeven is de lijn $l_1$ met de vergelijking $-3x = -y + 7$. Bepaal de vergelijking van de rechte lijn $l_2$, die gaat door het punt $Q(6,1)$ en die loodrecht staat op $l_1$. \\

\item[e.] De rechte lijn met vergelijking $2x + 3y = 4$ staat loodrecht op de lijn $y = ax + b$ die door het punt $A(1,2)$ gaat. Bepaal $a$ en $b$. \\
\end{enumerate}

\subsection*{Opgave 5: Snijpunten.}
\begin{enumerate}
\item[a.] Bereken het snijpunt van de lijnen $l_1$ en $l_2$ met de vergelijkingen $y = -x + 2$ en $y = x - 1$. \\
\item[b.] Bereken het snijpunt van de lijnen $l_1$ en $l_2$ met de vergelijkingen $y = 3x + 1$ en $y = -x + 9$. \\
\end{enumerate}